% Transcription — fonds Alexandre Grothendieck, shelfmark 27, batch 2 (pages 21-36).
% Established from the facsimile archives/batches/27/batch-02.pdf (cut from a
% copy of https://grothendieck.umontpellier.fr/27.pdf; PDF page k = archivists'
% page 20+k, no cover sheet), per the skill transcribe-grothendieck. The folder
% is thirty-six pages; this is the second and last batch (pages 1-20 are
% batch 1). The rules follow batch 1's header: the typed text is published and
% is not re-set; each typed page gets one French \note; typed versos with
% nothing of his are skipped.
%
% Pages skipped: 30 (below). Page 27 is a tan cover sheet inscribed in ink,
% top right and underlined, « Théorème de 1/2 continuité » (the 1/2 written as
% a small fraction, read as « semi- » per hand.md); it takes no \page marker,
% and his words become the section heading with a \note.
%
% Page 30 is also skipped: the other side of the sheet of page 29, scanned
% upside down, a typed leaf from a letter on ranks of modules of differentials,
% not on the subject of the folder and with nothing of his that can be told
% apart (a few symbols are inked into the typing's blanks, in a hand not
% identifiable as his) — a typed verso, skipped by the settled rule.
%
% Pages transcribed from his hand: 21-26, 28-29, 31-36. Page 22 was scanned
% upside down and was read turned; its whole body is cancelled by two long
% diagonals, and it is transcribed in full with the cancellation recorded once
% in a \note (hand.md §6). Page 26 likewise
% under three long strokes. Fast drafting throughout: the formulas and the
% numbered conditions and inequalities carry, much of the connecting prose does
% not, and is left \ill{}.
%
% His pagination: none on these leaves. Pages 24-25 are numbered 3) and
% (implicitly) the continuation of the numbered remarks 1), 2) of page 21;
% page 23 ends on « 3) » and page 24 opens the corresponding example. Pages
% 31-36 carry his own circled numbering (1)-(10) of the inequalities, fixed on
% page 29.
%
% What the batch is about. Two runs. (A) Pages 21-26, blue ink, the tail of
% the Lie-algebra material of batch 1 pp. 17-19 (complements to SGA 3 XIII on
% Cartan subalgebras and regular elements): p. 22 lists eleven conditions
% (i)-(xi) on a Lie subalgebra h of g (h contains a Cartan subalgebra, h = u
% and contains regular elements, ad(a) on g/h injective, psi: X -> W(g)
% dominant / generically smooth / quasi-finite, h the Lie algebra of a smooth
% subgroup, N smooth), with implication arrows and, below, a list a)-e) of
% candidate definitions (h equal to its normaliser and satisfying the density
% theorem, ...). P. 21 gives remarks with counterexamples in characteristic 2:
% SL(2,k), where g is nilpotent and psi is birational yet a condition fails;
% PGL(2,k) (« Pl(2,k) ») with basis X, Y, H, [X,Y] = 2H = 0, continued on
% p. 23 (h = kH + k(lambda X + mu Y) is not the Lie algebra of any smooth
% subgroup). P. 24: an example of a smooth solvable group G = T.U, T = G_m^2,
% U = G_a^I, over F_p, whose Lie algebra has no regular element (the kernels
% of the alpha_i cover F_p^2). P. 25: a group scheme over a DVR whose centre is
% not representable (E with trivial centre, F a nontrivial commutative subgroup,
% G obtained from E_S by removing E_k - F_k from the special fibre;
% Z(Â) = {e} while lim Z(A_n) = F). P. 26: a cancelled « Contrexemple ! »
% about regular elements and Cartan subalgebras. (B) Pages 27-36, black ink on
% tan paper, « Théorème de 1/2 continuité »: p. 28 defines, for G smooth
% connected over an algebraically closed field with 0 -> H -> G -> A -> 0,
% the numerical invariants rho_ab = dim A, rho_aff = dim H, rho_r (reductive
% rank), rho_s, d_s, rho_n (dimension of Cartan subgroups), rho_na, rho_u,
% rho_ss, rho_rad, with their relations; p. 29 is a sheet of scratch work
% fixing the list (1)-(10) and the implications between them; p. 31 states the
% theorem: for G a smooth separated group prescheme of finite presentation
% over S, rho_aff, d, rho_n, rho_na, rho_u are upper semicontinuous and
% rho_ab, rho_ab + rho_r, rho_s lower semicontinuous; pp. 31-36 give the
% proof by reduction to a DVR, a split maximal torus in the special fibre,
% prolongation of the n-torsion _nT_0 to finite étale subgroups H(n), their
% centralisers giving a Cartan subgroup, a lemma on nilpotent groups, and for
% the lower semicontinuity of rho_ab the action of pi_1 of a henselian trait
% on the Tate module M ≃ Z_l^(2 rho_ab + rho_r) through the cyclotomic
% character; p. 36 ends with (3) <=> (9) and the semisimple rank via
% « n-rang ».
%
% Notation fixed once for the batch: gothic g, h, u, m_a, l (for his script
% L of p. 21) as \mathfrak; W(g) as written; G_m, G_a bold as in batch 1; the
% invariants as \rho_{\mathrm{ab}}, \rho_{\mathrm{aff}}, \rho_{\mathrm{r}}
% etc., primes for the generic fibre as he writes them (rho', d'); his circled
% numbers as (1)-(10) in plain parentheses; s_0, s_1 the closed and generic
% points, bars for geometric points.
%
% Readings to check first: « Même si » for his « M̄ si » (p. 21); the condition
% numbers « (ix) », « (x) » on p. 21; the list a)-e) of p. 22; the whole
% prose of p. 24 (« caractères infinitésimaux », the margin); « Contrexemple »
% and the objects S, G of p. 26; « rho_as » against « rho_ab » at the top of
% p. 28 (kept as rho_ab); the margin « G lisse/k » (p. 33) and the Lemme at its
% foot; « hensélien strict. local » (p. 34); « caractère cyclotomique » and
% « n-rang » (p. 36).
%
% Pass: Opus 5.5 (claude-opus-5-5), 2026-09-24 — first pass, unchecked against the pages by a human.

\documentclass[11pt,a4paper]{article}
% Le préambule commun à toutes les transcriptions du fonds Grothendieck.
%
% Il tient en une idée : **l'incertitude se compose, elle ne se résout pas.**
% Une transcription qui lisse un mot douteux a détruit la seule chose pour
% laquelle on la faisait — un lecteur doit pouvoir savoir, à chaque endroit,
% ce qui était sur le feuillet et ce qui a été ajouté. Les six macros qui
% suivent sont donc l'appareil critique, et non de la décoration.
%
% Les mêmes commandes sont reconnues par `scripts/render.mjs`, qui produit la
% vue de lecture du site. Une macro ajoutée ici doit l'être aussi là-bas,
% faute de quoi le rendu échouera bruyamment — ce qui est voulu : un rendu
% silencieusement faux est pire qu'un rendu absent.

\ProvidesPackage{grothendieck}[2026/08/08 Transcriptions du fonds Grothendieck]

\usepackage{amsmath,amssymb,amsthm}
\usepackage{mathrsfs}
\usepackage{stmaryrd}
% Les diagrammes commutatifs sont partout dans ce fonds : c'est la langue
% même de la géométrie algébrique de Grothendieck, et une transcription qui
% les rendrait en prose ne transcrirait rien.
\usepackage{tikz-cd}
% Français par défaut — c'est la langue des feuillets — mais l'anglais est
% chargé aussi : la traduction et le résumé sont des documents anglais, et la
% typographie française (espaces avant les deux-points) y serait une faute.
% Ces documents basculent par \selectlanguage{english} en tête de corps.
\usepackage[english,french]{babel}
\usepackage{fontspec}
\usepackage{geometry}
\geometry{a4paper,margin=25mm,marginparwidth=28mm}
\usepackage{marginnote}
\usepackage{xcolor}
% ulem, not soul. `soul`'s \st and \ul are built on a character-by-character
% scan that XeLaTeX's font machinery breaks: the document fails with an
% undefined control sequence rather than degrading. ulem does the same two jobs
% and is Unicode-engine safe. `normalem` stops it hijacking \emph, which the
% transcriptions use for Grothendieck's own emphasis.
\usepackage[normalem]{ulem}
\usepackage{hyperref}
\hypersetup{colorlinks=true,linkcolor=black,urlcolor=black,pdfborder={0 0 0}}

% --- Métadonnées du lot -----------------------------------------------------
% Reprises telles quelles par le script de rendu, qui les affiche en tête de
% la vue de lecture. Elles fixent ce que le fichier prétend couvrir : sans
% elles, une transcription flotte sans point d'attache dans un fonds de seize
% mille feuillets.

\newcommand{\folder}[1]{\gdef\grothFolder{#1}}
\newcommand{\batch}[1]{\gdef\grothBatch{#1}}
\newcommand{\leaves}[2]{\gdef\grothFirst{#1}\gdef\grothLast{#2}}
\newcommand{\foldertitle}[1]{\gdef\grothFoldertitle{#1}}
\gdef\grothFolder{?}\gdef\grothBatch{?}\gdef\grothFirst{?}\gdef\grothLast{?}\gdef\grothFoldertitle{}

% --- Le résumé --------------------------------------------------------------
%
% Il ouvre la lecture modernisée, sous le titre « Résumé », et il est composé
% en petit. Ce n'est pas un ornement : c'est le seul passage du document qui
% ne soit pas des mathématiques, et un lecteur doit pouvoir le reconnaître
% avant de l'avoir lu — puis le sauter s'il connaît déjà le sujet, ou s'y
% arrêter s'il ne le connaît pas. Le filet qui le suit marque la reprise du
% corps.

\newenvironment{resume}
  {\small\setlength{\parskip}{0.45em}}
  {\par\medskip\hrule\medskip}

% --- Mots-clés ---------------------------------------------------------------
%
% En anglais, en clôture du résumé de la lecture modernisée : le vocabulaire
% moderne sous lequel chercher ce que les feuillets construisent. Le manifeste
% du site (`npm run manifest`) les extrait pour étiqueter la cote entière —
% cette ligne est la source unique des tags, il n'y a pas de fichier à côté.
\newcommand{\keywords}[1]{\par\smallskip\noindent
  {\footnotesize\textsc{Keywords} --- \textit{#1}}\par}

% --- L'appareil critique ----------------------------------------------------

% Le numéro de feuillet, en marge. C'est l'ancre qui permet de régler un
% désaccord sur une formule en regardant *un* feuillet plutôt qu'une chemise
% de six cents. Le site s'en sert aussi pour tourner les pages du fac-similé
% quand on fait défiler la transcription.
\newcommand{\leaf}[1]{%
  \marginnote{\footnotesize\color{gray!70}\textsf{#1}}%
  \label{leaf:#1}\ignorespaces}

% Illisible : on ne devine pas. Un mot inventé qui se lit comme les autres est
% le pire résultat possible.
\newcommand{\ill}{\textcolor{red!60!black}{[\,\ldots\,]}}

% Lecture douteuse : le mot est donné, et signalé comme incertain.
\newcommand{\uncertain}[1]{\uline{#1}}

% Ajout de l'éditeur — une lettre manquante, un mot sous-entendu.
\newcommand{\add}[1]{\textcolor{blue!50!black}{[#1]}}

% Biffé par Grothendieck lui-même. Ce qu'il a rayé fait partie du document :
% une rature dit souvent où la pensée a changé de direction.
\newcommand{\struck}[1]{\sout{#1}}

% Note du transcripteur, sur l'état matériel du feuillet.
\newcommand{\note}[1]{\par\smallskip{\small\color{gray!60!black}\textit{[#1]}}\par\smallskip}

% Note marginale de Grothendieck, distincte du corps du texte.
\newcommand{\marginal}[1]{\par\smallskip\noindent\hspace{1em}%
  {\small\color{gray!60!black}#1}\par\smallskip}

% --- Titre ------------------------------------------------------------------

\AtBeginDocument{%
  \begin{center}
    {\large\bfseries Fonds Alexandre Grothendieck --- cote n\textsuperscript{o}~\grothFolder}\\[2pt]
    {\itshape\grothFoldertitle}\\[4pt]
    {\small feuillets \grothFirst--\grothLast\ (lot \grothBatch)}\\[2pt]
    {\footnotesize\color{gray!60!black}Transcription établie d'après le fac-similé numérisé
     par l'Université de Montpellier.\\
     Les lectures incertaines sont soulignées, les illisibles notées [\,\ldots\,],
     les ajouts entre crochets.}
  \end{center}
  \vspace{1em}}

\endinput


\folder{27}
\batch{2}
\pages{21}{36}
\dating{[vers 1962-1970]}
\watermark{Édition de démonstration}
\foldertitle{SGA D [SGA 3]. Compléments XII à XVII : tapuscrit annoté (s.d.), notes manuscrites (s.d.)}

\begin{document}

\page{21}

\note{encre bleue ; suite, semble-t-il, des notes sur les algèbres de Lie des
pages 17 à 19 ; les numéros de conditions renvoient à la liste de la page 22}

\textbf{Remarques} \textbf{1)} \uncertain{Même si}
$\psi$ est \struck{\ill{}} \add{birationnel}, \emph{et}
$\mathfrak{h}$ provient d'un sous-groupe $H$ de $G$ lisse sur $k$,
\struck{\ill{}} la condition \uncertain{(ix)} \ill{} \uncertain{pas}
\uncertain{nécessairement} \uncertain{vérifiée}.\note{« Même si » est écrit
« M̄ si », ici et en 2) ; « birationnel », souligné, est en interligne
au-dessus d'un mot biffé, peut-être « étale »}

\emph{Exemple} \quad $G = \mathrm{Sl}(2,k)$, $k$ car.~$2$, alors
$\mathfrak{g}$ est \emph{nilpotente}, donc la \add{seule}
sous-algèbre de Cartan \struck{\ill{}} est $\mathfrak{g}$ elle-même. Soit
$B$ Borel, alors \struck{$\mathfrak{L}$} $\mathfrak{h} = \mathrm{Lie}(B)$,
on sait que
\[
  \bigcup_{g \in G(k)} \mathrm{ad}(g)\, \mathfrak{L} = \mathfrak{g}
\]
\struck{et} \ill{}, donc \struck{\ill{}} donc $\psi$ dominant. D'ailleurs
\[
  \mathfrak{u} = \mathfrak{g}, \quad B \subset N \neq G, \quad \text{donc}
  \quad \dim N = \dim B = \mathrm{rang}\, \mathfrak{h},
\]
donc $\dim X = \dim G$, \struck{D'ailleurs \ill{} \ill{}}
\add{donc $\psi$ \uncertain{gén.} \uncertain{qu.-fini}}.

\struck{\uncertain{avec} que $\mathfrak{u} = \mathfrak{e}_{\mathfrak{h}}$,
\quad $\mathfrak{u} = \mathfrak{m}_a = \mathfrak{g}$ (car $\mathfrak{u}
\subset \mathfrak{m}_a$)}

D'ailleurs, \struck{\ill{}} \ill{} plus \ill{} \ill{} \struck{donc $\varphi$
\emph{lisse} en tout point \ill{}, $\psi$ étale en tout point}\note{deux
lignes barrées et couvertes d'une ligne ondulée ; à droite, un encadré
entièrement raturé où se lisent encore $N = M_a$, $G/N$, $B \subset G$ et
deux points marqués $a$, $b$}

…\ ici que $\psi$ \emph{est birationnel}.
\struck{\ill{} \ill{} $B \cap B$}

\textbf{2)} \uncertain{Même si} $\psi$ \struck{birationnel}
\add{dominant} et $\mathfrak{h} = \mathfrak{u}$, la condition
\uncertain{(x)} \ill{} \uncertain{pas} \uncertain{nécessairement}
\uncertain{vérifiée}. Prenons $G = \mathrm{Pl}(2,k)$, $k$ car.~$2$,
$\mathfrak{g}$ \ill{} base $X, Y, H$,
\[
  [X, Y] = 2H = 0, \qquad [H, X] = X, \qquad [H, Y] = Y.
\]
\note{suite à la page 23}

\note{encadré à gauche, sous un trait horizontal fléché :}
\textbf{N.B.} $\varphi$ lisse en $(e, a)$
\[
  \Updownarrow
\]
$\mathrm{ad}(a)_{\mathfrak{g}/\mathfrak{h}}$ injectif, i.e.
\[
  [x, a] \in \mathfrak{h} \Rightarrow x \in \mathfrak{h}, \qquad
  \mathfrak{m}_a = \mathfrak{h}
\]
\[
  \Downarrow
\]
\[
  \mathfrak{h} = \mathfrak{u}
\]

\page{22}

\note{page écrite tête-bêche par rapport au scan, lue retournée ; tout le
corps est annulé par deux longues diagonales, et se lit néanmoins. Dans la
marge gauche, une accolade réunit (i)-(iii), des flèches doubles verticales
relient les conditions (v)-(ix), une accolade (ix)-(xi), et un long trait
courbe remonte de (xi) vers (iii)}

\marginal{$\mathfrak{h} \subset \mathfrak{u} \subset \mathfrak{m}_a$}\note{en
haut à droite, souligné d'un trait fléché}

\begin{itemize}
  \item[(i)] $\mathfrak{h}$ contient une sous-algèbre de Cartan
  \item[(ii)] $\mathfrak{h} = \mathfrak{u}$, $\mathfrak{h}$ contient
  \uncertain{des éléments} réguliers de $\mathfrak{g}$
  \item[(iii)] $\exists\, a \in \mathfrak{h}$ t.q.
  $\mathrm{ad}_{\mathfrak{g}/\mathfrak{h}}(a)$ \struck{i}injectif (i.e.
  tel que \struck{\ill{}} $\mathfrak{h} = \mathfrak{m}_a$)
  \item[(iv)] $\psi : X \to W(\mathfrak{g})$ dominant, et $\mathfrak{h} =
  \mathfrak{u}$.
  \item[(v)] $\psi : X \to W(\mathfrak{g})$ gén.\ lisse
  \item[(vi)] $\psi : X \times W(\mathfrak{h}) \to W(\mathfrak{g})$
  \uncertain{gént} lisse
  \item[(vii)] $\psi : X \to W(\mathfrak{g})$ dominant
  \item[(viii)] $\psi : X \to W(\mathfrak{g})$ dominant et quasi-fini
  \item[(ix)] $\psi$ gén.\ qu.-fini et $\mathfrak{h} = \mathfrak{u}$
  \item[(x)] \struck{(xiii)} \struck{$\psi$ dominant} $\psi : X \to
  W(\mathfrak{g})$ dominant, \emph{et} [$\mathfrak{h}$ est l'alg.\ de Lie
  d'un sous-groupe lisse $H$ de $\mathfrak{g}$]
  \item[(xi)] \struck{(xii)} $\psi : X \to W(\mathfrak{g})$ dominant, et $N$
  lisse.
\end{itemize}
\note{un point noir en marge droite en regard de (iv), (vii), (viii), (x),
(xi)}

$\mathfrak{h}$ \struck{\ill{}} \ill{} \uncertain{minimale}
\uncertain{parmi} les $k$-algèbres de Lie qui \ill{} :

\begin{itemize}
  \item[a)] soit $=$ \uncertain{normalisateur}, et \uncertain{satisfasse}
  \ill{} \add{\uncertain{densité}} ;
  \item[b)] soit $=$ \uncertain{normalisateur}, et \uncertain{satisfait} au
  th.\ de \uncertain{densité} \add{\uncertain{densité}} ;
  \item[c)] soit définie par $H$ lisse, et
  \item[d)] soit telles que $\exists\, a \in \mathfrak{h}$,
  $\mathrm{ad}(a)_{\mathfrak{g}/\mathfrak{h}}$ injectif
  \item[e)] \struck{soit telles que}
\end{itemize}
\note{un mot souligné et répété au-dessus de a) et sous b), lu
« densité » avec doute ; « th.\ de densité » est souligné d'un long trait}

\page{23}

\note{suite de l'exemple 2) de la page 21}

La sous-algèbre de Cartan est $\mathfrak{d} = \mathfrak{t} = kH$.
Ici, \struck{\ill{} \ill{} \ill{} \ill{} \ill{}} \ill{} de $\mathfrak{g}$
\struck{\uncertain{contenant}} $\mathfrak{t} = kH$ \struck{\ill{} \ill{} \ill{}}
\add{est une sous-algèbre}. \struck{\ill{} \ill{} \ill{} $=$ Cartan}

Mais prenons
\[
  \mathfrak{h} = kH + k(\lambda X + \mu Y),
\]
si \struck{$\mathfrak{h}$ était} \struck{\ill{}} le normalisateur
\add{$N/e$} de $\mathfrak{h}$ étant de dim~$2$, (donc $\mathfrak{h} =
\mathrm{Lie}(H)$) \struck{il contiendrait $T'$} i.e.\ lisse) donc dans
\struck{il contiendrait un tore maximal $T'$ \ill{} \ill{} \ill{}}
\struck{\ill{}} \ill{} \ill{}, \struck{\ill{} fait} \add{un Borel}, or
\struck{\ill{} \ill{}} \ill{} \ill{} (T est un \uncertain{C-sous-groupe})
\struck{\ill{}} $\mathfrak{h} \supset \mathfrak{t}$ \struck{\ill{} Borel
\ill{} \ill{} $B$, et tel que} \struck{$\mathfrak{h} = kH + kX$ \ill{}
\ill{} $B = \mathrm{Norm}\ldots$} \uncertain{impliquerait} $H \supset T$,
donc on \ill{} \ill{} \ill{} Borel(s)\note{lignes très raturées ; un encadré
barré d'une ondulation à gauche}

\struck{donc} \ill{} si $\lambda, \mu \neq 0$, \ill{} contenant $T$,
\struck{\ill{}} donc \ill{} $\mathfrak{h}$ n'est pas \uncertain{algèbre} de
Lie d'\uncertain{aucun} sous-groupe lisse de $G$ …

\textbf{3)}

\page{24}

\textbf{Exemple} d'un groupe lisse résoluble \add{\uncertain{connexe}}
$G$ sur $\mathbb{F}_p$, \add{\ill{} \ill{} \ill{}}\note{interligne entouré,
en haut à droite, en partie biffé} dont l'algèbre de Lie n'a pas d'élément
régulier. \qquad $\mathfrak{t} = \mathbb{F}_p^2$.

Soit $(\alpha_i)_{i \in I}$ une famille \add{\uncertain{finie}} de formes
\add{\uncertain{lin.} $\neq 0$} sur \struck{\ill{}} $\mathbb{F}_p^2$ telles
que
\[
  \bigcup \mathrm{Ker}\, \alpha_i = \mathbb{F}_p^2 .
\]
\note{le second membre a été récrit plusieurs fois ; la lecture
$\mathbb{F}_p^2$ est celle de la dernière couche}
(On peut \uncertain{prendre} \ill{} \ill{} \uncertain{formes} \ill{}
\uncertain{nulles}, et $\bigcap \mathrm{Ker}\, \alpha_i = 0$). Soient
\struck{\ill{}} les $\chi_i$ \ill{} \uncertain{caractères}
\uncertain{infinitésimaux} associés : \ill{} \uncertain{homomorphismes}
\[
  \chi_i : T = \mathbf{G}_m^2 \longrightarrow \mathbf{G}_m
\]
qui \ill{} \ill{} \ill{} \ill{} \ill{} \ill{}
\[
  G = T \cdot U \struck{\mathbf{G}_a^{I}}, \qquad U = \mathbf{G}_a^{I}
\]
avec $T$ opérant sur $\mathbf{G}_a^{I}$ par les $\chi_i$.

Alors $T$ opère sur $\mathfrak{g}/\mathfrak{t}$ par les
\uncertain{caractères} $\chi_i$, on peut \uncertain{montrer} que $T$
\ill{} \ill{} \ill{} i.e.\ le \uncertain{centralisateur} de $T$ \ill{} \ill{}.
Le centre de $G$ est donc \uncertain{égal} au centre \uncertain{réductif},
donc égal au \uncertain{noyau} de $T \to \mathrm{GL}(\mathfrak{g})$ ; qui
est \uncertain{nul}, \struck{\ill{} \ill{}} Enfin, un \struck{\ill{}}
élément de $\mathfrak{t}$ \struck{\ill{} \ill{}} étant dans
$\bigcup \mathrm{Ker}\, \alpha_i$, ne peut être régulier. Or \emph{tout}
\emph{tore maximal} \emph{de $G$} \emph{est} \emph{conjugué de $T$}
[\uncertain{par} un élément unique de $U(\mathbb{F}_p)$] donc
\struck{\ill{}} \struck{\ill{}} $\mathfrak{g}$ n'a pas d'élément
\uncertain{régulier} [car il serait conjugué d'un élément de
$\mathfrak{t}$].

\marginal{\struck{\ill{}} $T$ est \struck{\ill{}} \ill{} \ill{} \ill{}
i.e.\ $C(T)$ \ill{} \ill{} \ill{} \ill{} \ill{} \ill{}}\note{note écrite
de biais dans la marge gauche, en regard de « Le centre de $G$ … »}

\page{25}

\struck{Soit}

Cas d'un schéma en groupes \add{\uncertain{séparé}} étale de
\uncertain{t.f.} sur $S$ dont le centre n'est pas représentable
($S =$ spectre d'un anneau de valuation discrète \struck{$A$}~$A$).

On part d'un groupe ensembliste $E$, de centre trivial, et d'un
\uncertain{sous-groupe} commutatif \struck{\ill{}} $F$ non
\struck{\ill{}} \add{trivial} (p.ex.\ le groupe engendré par un élément
$s \neq e$ de $E$). Soit \struck{$\mathcal{F}$} $G$ le sous-groupe ouvert de
$E_S$ obtenu en enlevant de la fibre spéciale $E_k$ de $E_S$ le fermé
$E_k - F_k$. Donc
\[
  G_0 = F_k, \qquad G_1 = E_K .
\]
Alors le foncteur $\mathrm{Cent}(G)$ n'est pas représentable, car il ne
\uncertain{commute} pas : \ill{} \ill{} \uncertain{projective}
\uncertain{indiquée} des $A_n = A/\mathfrak{m}^{n+1}$, on a donc
\[
  Z(\hat{A}) = \{e\}, \qquad \varprojlim Z(A_n) = F .
\]

\textbf{N.B.} Si $G$ est \add{\uncertain{plat}} \struck{\ill{}} de
présentation finie sur $S$, à fibres connexes \uncertain{alors} $Z$ est
représentable ……

\page{26}

\note{la page est barrée de trois longs traits obliques}

\struck{\emph{\uncertain{Contrexemple} !}} Soit $\mathfrak{s}$ une
\uncertain{sous-algèbre} de $\mathfrak{g}$\note{les deux lettres sont
de forme cursive ; lecture douteuse}

\begin{itemize}
  \item[(i)] Si \struck{\ill{}} $\mathfrak{s}$ \ill{} tel que
  $\mathfrak{g}_{\mathfrak{s}}$ \emph{régulier}, \ill{} que
  \uncertain{s'annule} \ill{} $\mathfrak{s}$, $\mathfrak{g}_{\mathfrak{s}}$
  \ill{} \emph{régulier}.
  \item[(ii)] Si $\mathfrak{g}$ est régulier, alors $\exists$ une
  \uncertain{sous-algèbre} de Cartan \ill{} une \ill{} $\mathfrak{c}$ de
  $\mathfrak{g}$ qui contient $\mathfrak{s}$.
\end{itemize}

\struck{\emph{\uncertain{Contrexemple}}}\note{mot souligné, repassé puis
biffé ; le reste de la page est vide}

\section*{Théorème de ½ continuité}

\note{titre de sa main, souligné, en haut à droite du feuillet de garde
(page 27), qui ne porte rien d'autre ; le « ½ » est écrit en petite fraction,
pour « semi- ». Le feuillet ne reçoit pas de numéro de page}

\page{28}

$G$ lisse connexe sur le corps alg.\ clos.
\[
  e \longrightarrow H \longrightarrow G \longrightarrow A \longrightarrow e
\]
\note{sous $H$ : « groupe affine lisse connexe » ; sous $A$ : « abélienne »}

\[
  \rho_{\mathrm{ab}}(G) = \dim A \qquad \rho_{\mathrm{aff}}(G) = \dim H
\]
\note{le premier indice se lit aussi « as »}
\[
  \rho_{\mathrm{r}}(G) = \rho_{\mathrm{r}}(H) = \text{dim tores maximaux}
\]
\[
  \rho_{\mathrm{s}}(G) = \rho_{\mathrm{s}}(H) =
  \rho_{\mathrm{r}}(H/\mathrm{rad}(H))
\]
\note{avant $\rho_{\mathrm{r}}$, deux signes biffés, le second « dim »}
\[
  d_{\mathrm{s}}(G) = d_{\mathrm{s}}(H) = \dim(H/\mathrm{rad}(H)),
\]
\[
  \rho_{\mathrm{n}}(G) = \rho_{\mathrm{u}}(V) + \rho_{\mathrm{ab}}(G) =
  \text{dimension des Cartan}
\]
\[
  \rho_{\mathrm{na}}(G) = \rho_{\mathrm{u}}(V)
\]
\[
  \rho_{\mathrm{u}}(G) = \rho_{\mathrm{u}}(V) = \rho_{\mathrm{na}}(G) -
  \rho_{\mathrm{r}}(H)
\]

\note{à droite, une accolade :}
\[
  \left\{
  \begin{array}{l}
    \rho_{\mathrm{rs}}(G) = \rho_{\mathrm{rs}}(H) =
      \rho_{\mathrm{r}}(H/\mathrm{rad}(H)) \\[2pt]
    \rho_{\mathrm{rad}}(G) = \rho_{\mathrm{r}}(\mathrm{rad}(H))
  \end{array}
  \right.
\]
\note{entre les deux lignes, une ligne biffée « $\rho_{\mathrm{rad}}(G) =
\rho_{\mathrm{rad}}(H) = \ldots$ » ; dans la dernière ligne, $\rho_{\mathrm{r}}(\mathrm{rad}(H))$ est écrit
au-dessus d'un $\rho_{\mathrm{r}}$ ; lecture d'ensemble douteuse}

\struck{\ill{}}, $\underline{\rho_{\mathrm{ab}}}$, $\rho_{\mathrm{r}}$,
\struck{$\rho_{\mathrm{s}}$}, $\rho_{\mathrm{na}}$, $\rho_{\mathrm{ss}}$,
$\rho_{\mathrm{rad}}$\note{$\rho_{\mathrm{r}}$ … $\rho_{\mathrm{na}}$ et
$\rho_{\mathrm{ss}}$, $\rho_{\mathrm{rad}}$ sont enfermés dans deux cadres}
\qquad \struck{$\dim G = \rho_{\mathrm{aff}} + \rho_{\mathrm{ab}} =
\rho_{\mathrm{rad}} + \rho_{\mathrm{ss}} + \rho_{\mathrm{ab}}$}

\marginal{$\underline{\rho_{\mathrm{aff}}}$, $\rho_{\mathrm{rad}}$}

\[
  \rho_{\mathrm{na}} = \rho_{\mathrm{r}} + \rho_{\mathrm{u}}
\]
\[
  \rho_{\mathrm{n}} = \rho_{\mathrm{na}} + \rho_{\mathrm{ab}} =
  \rho_{\mathrm{r}} + \rho_{\mathrm{u}} + \rho_{\mathrm{ab}}
\]

\struck{$\rho_{\mathrm{s}} \geqslant \rho_{\mathrm{rs}}$}\qquad
\struck{$\rho_{\mathrm{r}} = \rho_{\mathrm{rs}} + \rho_{\mathrm{rr}}$}
\[
  \left\{
  \begin{array}{l}
    \rho_{\mathrm{aff}} = \rho_{\mathrm{ss}} + \rho_{\mathrm{rad}} \\
    \dim G = \rho_{\mathrm{ab}} + \rho_{\mathrm{aff}} = \rho_{\mathrm{ab}} +
      \rho_{\mathrm{ss}} + \rho_{\mathrm{rad}}
  \end{array}
  \right.
\]

$\rho_{\mathrm{r}}, \rho_{\mathrm{u}}, \rho_{\mathrm{ab}}, \rho_{\mathrm{na}},
\rho_{\mathrm{n}}$ sont les mêmes pour $G$ et \uncertain{pour} un
ss-groupe \uncertain{de} Cartan …\note{un trait vertical dans la marge
embrasse ces deux lignes}

\page{29}

\note{feuille de calculs épars, sans ordre ; la moitié supérieure est
traversée de grands traits obliques qui semblent délimiter des zones plutôt
qu'annuler. On transcrit de haut en bas}

\[
  {}_n A_0 \qquad T_0 \subset C_0 \qquad
  T_0 \subset \underbrace{C_0^{H}}_{\text{\scriptsize V.A}} \subset C_0
\]
\note{le $C_0^{H}$ porte un petit $\circ$ en indice ; « V.A » sous
l'accolade}
\[
  (x u) \qquad x^{n} = 1 \qquad (xu)^{n} = \,??
\]
\[
  H^2(\ , T_0) \simeq H^2(\ , {}_\infty T_0) = H^2(\ , \mathbb{Q}/\mathbb{Z})
\]
\[
  G - \qquad \rho_{\mathrm{ab}} \leqslant \rho'_{\mathrm{ab}}
\]
\[
  \begin{array}{c}
    \rho_{\mathrm{ab}} + \ldots \\
    \vee\!| \\
    \rho'_{\mathrm{ab}}
  \end{array}
  \qquad \frac{G_1}{Z_1} = Z \qquad
  \rho_{\mathrm{ab}}(Z_0) \leqslant \rho_{\mathrm{ab}}(Z_1) =
  \rho_{\mathrm{ab}}(G_1)
\]
\note{à gauche, le « $+\,\ldots$ » après $\rho_{\mathrm{ab}}$ est
raturé}
\[
  \rho_{\mathrm{ab}} + 2\rho_{\mathrm{r}} \qquad
  \rho_{\mathrm{a}} + \qquad \rho_{\mathrm{ab}} + \qquad
  \rho_{\mathrm{rig}}.
\]
\note{deux traits verticaux sous « $G -$ » ; un long trait horizontal sépare
la moitié supérieure}

\note{suit un tableau d'accolades, qu'on rend ainsi : sous
$\rho_{\mathrm{naff}}$ s'accolent $\rho_{\mathrm{rs}}$, $\rho_{\mathrm{u}}$,
$\rho_{\mathrm{ab}}$, et sous eux $\rho_{\mathrm{n}}$ ; sous $d$ s'accolent
$d_{\mathrm{ss}}$ et $d_{\mathrm{radaff}}$, puis $d_{\mathrm{rad}}$ et
$d_{\mathrm{aff}}$ ; à gauche $\rho_{\mathrm{ss}}$}
\struck{$\rho_{\mathrm{r}} = \rho_{\mathrm{ss}} + \rho_{\mathrm{rad}}\,
\rho_{\mathrm{r.radaff}}$}
\[
  d_{\mathrm{radaff}} \geqslant \rho_{\mathrm{u}} + \tilde\rho_{\mathrm{r}}
\]
\[
  \rho_{\mathrm{ab}} + d_{\mathrm{ss}} + \underbrace{(\rho_{\mathrm{r}} -
  \rho_{\mathrm{ss}})}_{\rho_{\mathrm{r\,radaff}}} \qquad
  \rho_{\mathrm{ab}} + \rho_{\mathrm{r}} + (d_{\mathrm{ss}} -
  \rho_{\mathrm{ss}})
\]
\note{les deux expressions sont enfermées dans un cadre arrondi}

\note{à droite, deux colonnes numérotées, les nombres entourés ; la
première réunie par une accolade, la seconde par une autre}
\[
  \begin{array}{ll}
    (1)\ \rho_{\mathrm{ab}} & (5)\ \rho_{\mathrm{n}} \\
    (2)\ \rho_{\mathrm{r}} + \rho_{\mathrm{ab}} = \rho_{\mathrm{rig}} &
      (6)\ \rho_{\mathrm{u}} \\
    (3)\ \rho_{\mathrm{ss}} & (7)\ \rho_{\mathrm{naff}} \\
    (4)\ d_{\mathrm{ss}} & (8)\ d_{\mathrm{aff}} \\
     & (9)\ d_{\mathrm{rad}} \\
     & (10)\ d_{\mathrm{radaff}}
  \end{array}
\]
\[
  \begin{array}{l}
    (2) + (5) \Longrightarrow (6) \\
    (1) + (5) \Longrightarrow (7) \\
    (1) \Longleftrightarrow (8) \\
    (3) \Longleftrightarrow (9) \\
    (1) + (3) \Longrightarrow (10)
  \end{array}
\]
\note{le premier chiffre de la première ligne est surchargé (un 2 sur un
autre chiffre)}

\page{31}

\emph{Th} \quad Soit $G$ un préschéma en groupes \add{lisse} séparé de
présentation finie sur $S$. Alors

\begin{itemize}
  \item[(i)] \struck{Les fonctions} $\rho_{\mathrm{aff}}$,
  $d$\struck{$\dim G$} sont $\tfrac12$ continues sup., et de
  \struck{\ill{}} les fonctions $\rho_{\mathrm{n}}$, $\rho_{\mathrm{na}}$,
  $\rho_{\mathrm{u}}$ sont semi-continues supérieurement \struck{si $G$ est
  \ill{} sur $S$ [\ill{} \ill{} \ill{} \ill{}]}\note{« $\rho_{\mathrm{aff}}$,
  $d$ » entouré, en interligne, relié par un trait au début de la ligne
  suivante ; le « $d$ » est écrit sur un mot biffé}
  \item[(ii)] \struck{Si $G$ est \emph{plat} sur $S$}, les fonctions
  $\rho_{\mathrm{ab}}$, $\rho_{\mathrm{ab}} + \rho_{\mathrm{r}}$,
  $\rho_{\mathrm{s}}$ sont semi-continues inférieurement.
\end{itemize}
\note{un trait vertical dans la marge gauche embrasse l'énoncé}

On se ramène comme d'habitude au cas où $S$ \uncertain{affine}
\uncertain{noeth.}, spectre d'un anneau de valuation discrète, \ill{}
\uncertain{point} \uncertain{générique} $s_1$, et à \uncertain{prouver} les
inégalités
\[
  \left\{
  \begin{array}{l}
    \rho_{\mathrm{aff}} \geqslant \rho'_{\mathrm{aff}} \qquad
      d \geqslant d' \\[2pt]
    (5)\ \rho_{\mathrm{n}} \geqslant \rho'_{\mathrm{n}}, \quad
    (6)\ \rho_{\mathrm{u}} \geqslant \rho'_{\mathrm{u}}, \quad
    (7)\ \rho_{\mathrm{na}} \geqslant \rho'_{\mathrm{na}} \\[4pt]
    (1)\ \rho_{\mathrm{ab}} \leqslant \rho'_{\mathrm{ab}}, \quad
    (2)\ \rho_{\mathrm{ab}} + \rho_{\mathrm{r}} \leqslant \rho'_{\mathrm{ab}}
      + \rho'_{\mathrm{r}}, \quad
    (3)\ \rho_{\mathrm{s}} \leqslant \rho'_{\mathrm{s}}
  \end{array}
  \right.
\]
\note{numéros entourés ; le (5) et le (6) sont surchargés (des numéros
antérieurs corrigés), le (7) aussi ; « $\rho_{\mathrm{aff}} \geqslant
\rho'_{\mathrm{aff}}$ » est biffé en fin de la deuxième ligne et reporté,
entouré, en tête de la première, avec « $d \geqslant d'$ »}

\struck{\ill{}} \add{\uncertain{quitte}}\note{« quitte » en interligne,
sur un mot biffé} \uncertain{prenons} \ill{} \uncertain{réserve} de
\uncertain{platitude}. Ce sont \struck{\ill{}} \ill{} \uncertain{pour}
un groupe algébrique sur un corps \ill{}, \ill{} des \uncertain{quantités}
invariantes \ill{} \uncertain{changeant} $G$ en $G_{\mathrm{réd}}^{\circ}$,
et de plus pour $\rho_{\mathrm{aff}}$ et $d$, c'est une fonction
\emph{croissante} du groupe. Donc \ill{} \ill{} \uncertain{prouver}
[quitte à \uncertain{remplacer} $G$ par $\overline{G_1^{\circ}}$]
\struck{de \uncertain{remplacer} $G_1$ \ill{}} \ill{} \ill{} \ill{}
\emph{\uncertain{toujours}} $G$ plat sur $S$, et de plus

\page{32}

(quitte à faire une \emph{extension} finie de $S$ et
\uncertain{remplacer} $G$ par $\overline{G_{1\,\mathrm{réd}}^{\circ}}$\,]
que $G_1$ est \add{connexe et} lisse$/k(s_1)$, enfin [quitte à remplacer
$G$ par $G_0$ \uncertain{composante} neutre de $G_0$] que $G_0$ est
\uncertain{également} connexe [mais peut-être pas lisse].
\note{le « [ » ouvrant du premier crochet est à la fin de la page 31}

\textbf{a)} L'inégalité $d \geqslant d'$ bien connue, mais pour mémoire,
elle se réduit \emph{ici} (en effet) à l'égalité $d = d'$.
L'inégalité $\rho_{\mathrm{aff}} \geqslant \rho'_{\mathrm{aff}}$
\emph{équivaut alors} à $\rho_{\mathrm{ab}} \leqslant
\rho'_{\mathrm{ab}}$ [car $\rho_{\mathrm{aff}} + \rho_{\mathrm{ab}} =
\rho'_{\mathrm{aff}} + \rho'_{\mathrm{ab}} = d$], \uncertain{que} nous
allons traiter plus bas.

\textbf{b)} \struck{Considérons} On peut \uncertain{supposer} que $G_0$
admet un tore maximal splitté $T_0$ (car \ill{} \ill{} \ill{} \ill{}
\ill{} \ill{} alg.\ clos]. \struck{\ill{} ${}_nT_0$ \ill{} Soit} $C_0$ le
Cartan associé. \uncertain{Pour} \ill{} \ill{} \add{$T_0$ \ill{}}
\struck{\ill{} \ill{} prolonger le premier \ill{} \ill{} \ill{} \ill{}
résiduelle.}

\marginal{$G$ lisse$/S$}

\struck{Alors pour les premiers \ill{} \ill{} \ill{} $n$,} \add{\ill{} à
la caractéristique résiduelle, considérons}
\struck{${}_nC_0 = \mathrm{Ker}\,[\,n^{\text{ième}} \text{ des } C_0\,]$.
\uncertain{Comme} $C_0/T_0 = M_0$ \ill{}}\note{la ligne
« ${}_nC_0 = \ldots$ » est barrée d'un long trait qui descend vers la
droite ; « $= M_0$ » est écrit sur un mot raturé}

…\ \uncertain{est} une extension d'une V.A.\ $A_0$ par un groupe
\uncertain{unipotent} \add{lisse} \uncertain{connexe} $U_0$, on voit
\uncertain{que} les extensions correspondantes de ${}_nA_0$ par $U_0$
splittent \uncertain{canoniquement}, \struck{\ill{} \ill{} \ill{} des
extensions} \ill{} \ill{} ${}_nU_0$, \ill{}\note{les trois dernières lignes
sont traversées de trois courbes obliques qui semblent les annuler}

\page{33}

${}_nT_0$, on peut le prolonger en un sous-groupe \add{fini} étale $H(n)$
de $G$, tel que $H(m) = {}_mH(n)$ si $m \mid n$. La famille des sous-groupes
fermés lisses $\underline{\mathrm{Cent}}\, H(n)$ de $G$ est stationnaire, sa
valeur limite $C$ est telle que
\[
  D_0 = \mathrm{Cent}_{G_0}(T_0) = C_0 .
\]
\marginal{$G$ lisse$/k$}
\struck{De} plus, on voit que $C_1$ \ill{} \struck{$=$} \ill{}
\uncertain{réductif} et \uncertain{nilpotent} que $C_1$ [i.e.\
\struck{\ill{}} $C_{\bar s_1}$ contient un Cartan de $G_{\bar s_1}$], c'est
un \struck{\ill{} \ill{}} \ill{} à \uncertain{vérifier}\note{« c'est un
… à vérifier » souligné d'une ligne ondulée}.

\marginal{$\ast$}
Il s'ensuit que \struck{$\rho_{\mathrm{u}}(G_1) = \rho_{\mathrm{u}}(C_1)$,
$\rho_{\mathrm{r}}(G_1) = \rho_{\mathrm{r}}(C_1)$}
$\rho_{\mathrm{r}}, \rho_{\mathrm{u}}, \rho_{\mathrm{ab}}, \rho_{\mathrm{n}},
\rho_{\mathrm{na}}$ sont les mêmes pour $G_1$ et $C_1$
\[
  \begin{array}{l}
    \rho_{\mathrm{u}}(G_1) = \rho_{\mathrm{u}}(C_1) \\
    \rho_{\mathrm{ab}}(G_1) = \rho_{\mathrm{ab}}(C_1) \\
    \rho_{\mathrm{n}}(G_1) = \rho_{\mathrm{n}}(C_1)
  \end{array}
\]
\note{au-dessus, une première ligne biffée identique à la suivante ; les
quatre lignes sont encadrées et traversées de deux traits
obliques ; les lettres des deux dernières sont douteuses}
et \uncertain{sont} \ill{} \ill{} \ill{} les mêmes pour $C_0$ et $G_0$,
\uncertain{donc} pour la formule (5)\note{le chiffre avant la parenthèse
se lit « 1) » ou « 5) »}, \ill{} est ramené au cas où
\[
  G = C,
\]
i.e.\ $G_0$ \struck{\ill{}} \ill{} \uncertain{nilpotent}. Alors la relation
\[
  \boxed{\rho_{\mathrm{n}} \geqslant \rho'_{\mathrm{n}}}
\]
triviale, car $\rho_{\mathrm{n}} = d = d'$.

\emph{Lemme} $C$ groupe nilpotent \uncertain{connexe} lisse sur corps
alg.\ clos $k$. Alors
\[
  \overbrace{C/\underline{\mathrm{Cent}}(C)^{\circ}_{\mathrm{réd}}}^{Z}
\]
est unipotent. Pour \ill{} $H$ \uncertain{groupe} de type multiplicatif
$H$,
\[
  \underline{\mathrm{Hom}}_{k\text{-gr}}(H, C) \simeq
  \underline{\mathrm{Hom}}_{k\text{-gr}}(H, Z).
\]
En \uncertain{particulier}, \ill{} \ill{} \ill{} : \ill{} \ill{}
${}_nC \simeq {}_nZ$ \ill{} \ill{} \ill{} de $C$, \ill{} \ill{} \ill{}

\marginal{\uncertain{pour} \ill{} \ill{} \ill{} \ill{}}\note{écrit de
biais dans la marge gauche, en bas}

\page{34}

\struck{\emph{Corollaire} Soit $G$ lisse sur $S$, et supposons que}
\struck{On sait que \ill{} tel ss-groupe \ill{} $S$}

Appliquons ceci au cas où $C = G_0$, alors si $S$ est
\uncertain{hensélien} \uncertain{strict.}\ local, les ${}_nC$ se
\uncertain{prolongent} de \emph{façon} \emph{unique} en des
sous-groupes $H(n)$ \struck{\ill{}} étales finis de $G$,
\uncertain{évidemment} centraux, et
\[
  {}_mH(n) = H(m) \quad \text{si } m \mid n .
\]

Notons que $\rho_{\mathrm{ab}}$ et $\rho_{\mathrm{r}}$, donc
$\rho_{\mathrm{ab}} + \rho_{\mathrm{r}} = \rho_{\mathrm{rig}}$ sont
fonctions croissantes d'un groupe \struck{lisse} alg.\
\add{[qu'on peut supposer lisse$/k(s_1)$ !]}

Soit $K$ \struck{\ill{}} l'adhérence de \add{la limite des $H(n)_1$}, et
$\overline{K}_1 = K_{1}^{\circ}$, \add{alors} \struck{\ill{} \ill{}}
\[
  \begin{array}{l}
    \rho_{\mathrm{ab}}(K_1) \leqslant \rho_{\mathrm{ab}}(G_1) =
      \rho'_{\mathrm{ab}} \\
    \rho_{\mathrm{r}}(K_1) \leqslant \rho_{\mathrm{r}}(G_1) =
      \rho'_{\mathrm{r}} \\
    \rho_{\mathrm{rig}}(K_1) \leqslant \rho_{\mathrm{rig}}(G_1) =
      \rho'_{\mathrm{rig}}.
  \end{array}
\]
Donc pour prouver (1) et (2), on est ramené à prouver la même chose pour
$K$ remplaçant $G$\add{, \uncertain{puisque} $K_0 \supset H(n)_0 =
{}_nG_0$} \uncertain{ceci} \uncertain{facilement} (\uncertain{Cartan}
…),\note{une flèche courbe, dans la marge gauche, relie « D'autre part »
à « Donc pour prouver »}

D'autre part, \ill{}
\[
  \rho_{\mathrm{ab}}(G_0) = \rho_{\mathrm{ab}}(K_0), \qquad
  \rho_{\mathrm{r}}(G_0) = \rho_{\mathrm{r}}(K_0)
\]

\ill{} \ill{} \ill{} la limite de $G$, mais on a gagné la
\emph{commutativité}, qui implique d'ailleurs $\rho_{\mathrm{n}} =
\struck{\ill{}}\, \rho'_{\mathrm{n}} = d$. L'inégalité (2)
\uncertain{équivaut} \uncertain{alors} à l'inégalité (6), \struck{donc}
\add{et} (1) à (7) et (8).\note{« (1) à (7) et (8) » : chiffres
entourés, sous la ligne}

\page{35}

est \uncertain{ramené} à prouver $\rho_{\mathrm{u}} \geqslant
\rho'_{\mathrm{u}}$, \add{puis $\rho_{\mathrm{aff}} \geqslant
d'_{\mathrm{aff}}$}\note{interligne relié par une accolade ; l'indice
se lit mal} \uncertain{Remplaçant} $G$ par l'adhérence de $G_1$
\uncertain{unipot.} [\uncertain{quitte} \ill{} \ill{}
défini$/k(s_1)$] on est ramené \add{pour $\rho_{\mathrm{u}} \geqslant
\rho'_{\mathrm{u}}$} à \uncertain{prouver} que si $G_1$ est
\emph{unipotent}, \emph{il en est de même de} $G_0$.

\marginal{[\uncertain{sans} \uncertain{restriction} \ill{}
\uncertain{limitation}]}\note{en marge gauche, en regard ; la même note
revient plus bas}

\uncertain{Facile} pour les points d'ordre fini.

Pour la deuxième assertion, on est ramené à prouver que si $G_1$ est
\emph{affine}, $G_0$ \uncertain{est} \emph{affine}.

\marginal{[\uncertain{sans} \uncertain{restriction} \ill{}
\uncertain{limitation}]}

On est ramené au cas où \struck{\ill{} \ill{}} $S$ \add{\uncertain{semi}
\uncertain{local} $\dim S = 1$}, \struck{\ill{}} $\mathbb{Z}$, \ill{}
schéma \add{intègre} de type fini \struck{sur}, (\uncertain{à}
\uncertain{priori} …). \struck{\uncertain{Considérons} \ill{} $G$ et
la \ill{} \ill{}} \ill{} \ill{} \uncertain{supposer} \emph{le tore
maximal de $G_1$ est splitté.} Donc \ill{} \ill{} $k(s_0)$ est
\uncertain{fini}. \uncertain{Passant} à la clôture
\uncertain{hensélienne} (mais pas strictement hensélienne !) en $s_0$,
d'où un $\tilde S$, $\tilde G$ sur $\tilde S$, \struck{\ill{} ${}_n\tilde
G$} soit $\ell$ distinct de la caractéristique résiduelle, et un
\uncertain{premier} \ill{} \ill{} sur $\tilde S$, \uncertain{sous}
$\tilde H(n)$. ${}_n\tilde G$ \ill{} \ill{} \ill{} \ill{} \ill{}.

\note{dans la marge gauche, un croquis : une flèche descendant vers « $S$ »,
un angle droit marqué d'un point, et un segment portant un point}

Alors $\pi_1(\tilde S) = \pi_1(s_0) \simeq \mathbb{Z}$ \uncertain{opère}
sur $H(n)_{\bar s_0} = M_n$, d'où, à la limite, des opérations de
$\mathbb{Z}$ sur
\[
  \varprojlim M_n = M \simeq \mathbb{Z}_\ell^{\,2\rho_{\mathrm{ab}} +
  \rho_{\mathrm{r}}}.
\]
D'autre part, \struck{\ill{} \ill{}, \uncertain{donc}} comme
\[
  M_\infty \simeq \varprojlim H(n)_{\bar s_1}
\]
\note{un $H(n)$ entouré, au-dessus d'un indice raturé}
[\struck{\ill{}} \uncertain{opération} déterminée mod une opération de
$\pi_1(\tilde S)$ sur $M_\infty$] d'où
\[
  M_\infty \subset \varprojlim_n ({}_nG)_{\bar s_1} =
\]
\struck{donc}

\page{36}

\[
  \varprojlim ({}_nT)_{\bar s_1} \simeq ({}_n\mathbf{G}_m)^{r}_{\bar s_1},
\]
donc les opérations de $\pi_1(\tilde S, \bar s_0)$ sur $M_\infty$ se font
par l'intermédiaire du « \uncertain{caractère}
\uncertain{cyclotomique} » l'opération \uncertain{sur} les racines de
l'unité. Or si on a $\rho_{\mathrm{ab}} \neq 0$, alors les opérations
\ill{} \ill{} \ill{} Galois
\[
  \pi_1(\tilde S, \bar s_0) = \pi_1(s_0, \bar s_0)
\]
\ill{} \ill{} \ill{} \uncertain{par} l'intermédiaire
\uncertain{exclusif} de \ill{} \uncertain{caractère} \ill{}.
On est ramené à \ill{} \ill{} \uncertain{sur} une V.A.\ \ill{} sur un
corps $\mathfrak{k}$ de type fini \ill{} \ill{} \uncertain{mais}
\uncertain{ici} \uncertain{suffisant} …]\note{le « [ » ouvrant manque
sur la page}

\note{un trait horizontal sépare la suite}

\marginal{\ill{} \ill{} \ill{} …}\note{écrit de biais dans la marge
gauche, devant un trait vertical qui embrasse les lignes suivantes}

Restent à prouver (3) et (9).

\[
  (3) \Longleftrightarrow (9),
\]
et \struck{\ill{}} \uncertain{résultat} \ill{} \ill{} \uncertain{que}
\emph{un} \emph{\uncertain{groupe} réductif} \ill{}
\emph{\uncertain{spécialise} en un réductif}\note{chiffres
entourés}

\textbf{(4)} On regarde un tore maximal $T_0$ de $G_0$, on prolonge
les ${}_nT$ \ill{} \ill{} $H(n)$ \ill{} \ill{} \uncertain{sous-groupes}
compatibles \ill{} \ill{} \ill{} \ill{} $G_1$
\[
  H(n)_1/R_1 \cap H(n)_1 , \quad \ldots
\]
\[
  r_0 = \rho_{\mathrm{ss}}(G_0) \leqslant \dim T_0/T_0 \cap R_0
\]
\ill{} de « \uncertain{$n$-rang} » $\geqslant \dim T_0/T_0 \cap R_0
\geqslant r_0$. \ill{} \ill{} \ill{} \uncertain{$n$-rang} $\leqslant
\rho'_{\mathrm{ss}}$, d'où $\rho'_{\mathrm{ss}} \geqslant
\rho_{\mathrm{ss}} \ldots$\note{« $\leqslant \dim T_0/T_0\cap R_0$ » est
en interligne, relié par une accolade à « $r_0 = \rho_{\mathrm{ss}}(G_0)$ »}

\note{dans la marge gauche, en regard de (4) :}
\[
  \begin{array}{cc}
    G & G_1 \\
    \cup & \cup \\
    \overline{R}_1 = R & R_1
  \end{array}
  \qquad G_0/R_0
\]

\end{document}
